\documentclass[11pt]{article}

\usepackage[T1]{fontenc}
\usepackage{amsmath,amssymb,amsthm}
\usepackage[margin=1in]{geometry}
\usepackage{microtype}
\usepackage[hidelinks]{hyperref}

\newtheorem{theorem}{Theorem}
\newtheorem{lemma}[theorem]{Lemma}
\title{Rigidity at the Ideal Endpoint for Equiangular Hyperbolic $4$-Cubes}
\author{Xiaojian Zeng\\
PBC School of Finance, Tsinghua University, 100190\\
\texttt{xiaojianzeng@foxmail.com}}
\hypersetup{
  pdftitle={Rigidity at the Ideal Endpoint for Equiangular Hyperbolic 4-Cubes},
  pdfauthor={Xiaojian Zeng},
  pdfsubject={Hyperbolic 4-polytope rigidity at the ideal endpoint}
}
\date{}

\begin{document}
\maketitle

\begin{abstract}
Let $\alpha _0=\arccos(1/3)$.  We prove that every finite-volume
convex hyperbolic $4$-polytope with the face lattice of the $4$-cube,
with no hyperideal or truncated vertices, and with all dihedral angles
equal to $\alpha _0$, is isometric to the regular ideal equiangular
$4$-cube.  The proof is a Lorentz--Gram argument.  The ideal vertex
rays make the four differences of opposite-facet normals spacelike;
their positive span forces all four previously undetermined opposite
Gram entries to equal $-5/3$.  A relation-compatible Lorentz isometry
then carries the eight oriented halfspaces to the cube
$[-1/2,1/2]^4$ in the Klein ball.  This addresses only the ideal
endpoint for cube combinatorics and makes no volume, higher-dimensional,
or priority claim.
\end{abstract}

\section{Statement and source comparison}

Egorov studies finite-volume acute equiangular hyperbolic
$4$-polytopes and introduces
\[
  \alpha _0=\arccos(1/3).
\]
At this angle every vertex is ideal, and a regular ideal cube is
constructed in the Klein model.  For cube combinatorics, the uniqueness
statement in Proposition~2.3 of \cite{Egorov2026} is restricted to
$\alpha>\alpha _0$; its proof expressly says that uniqueness at
$\alpha=\alpha _0$ is not needed.  The same restriction occurs in its
Proposition~5.1 and Remark~5.2.  The compact uniqueness step cited there
is traced to Andreev's rigidity work \cite{Andreev1970}.  The argument
below is instead an endpoint argument and does not use compactness.

We use ``face lattice of the $4$-cube'' literally: the eight facets can
be labelled in four specified opposite pairs $F_i^+,F_i^-$, and a sign
vector in $\{+,-\}^4$ specifies one of sixteen distinct vertices.  All
vertices are assumed to be finite or ideal; hyperideal vertices and
polar truncations are not included.

\begin{theorem}\label{thm:main}
Let $P\subset\mathbb H^4$ be a full-dimensional finite-volume convex
hyperbolic polytope with the face lattice of the $4$-cube.  Suppose
that its vertices are finite or ideal and that every dihedral angle is
$\alpha _0=\arccos(1/3)$.  Then $P$ is isometric to the regular ideal
equiangular cube $C^4_{\alpha _0}$.  In Klein coordinates its closure
is $[-1/2,1/2]^4$.
\end{theorem}

\section{Ideal vertex rays}

Work in $V=\mathbb R^{4,1}$ with Lorentzian product
\[
 \langle x,y\rangle=-x_0y_0+x_1y_1+\cdots+x_4y_4,
 \qquad
 \mathbb H^4=\{x:\langle x,x\rangle=-1,\ x_0>0\}.
\]
Choose outward unit spacelike normals $u_i^\pm$ to the facets
$F_i^\pm$.  Thus
\[
 P=\{z\in\mathbb H^4:\langle z,u_i^\pm\rangle\leq0
       \text{ for every }i,\pm\}.
\]
Facets belonging to different opposite pairs meet.  The outward-normal
angle formula, in the conventions also used in \cite[Section~2.1]{Egorov2026},
therefore gives
\begin{equation}\label{eq:cross}
 \langle u_i^\varepsilon,u_j^\delta\rangle=-\cos\alpha _0=-\frac13
 \qquad (i\ne j).
\end{equation}

For $\sigma=(\sigma_1,\ldots,\sigma_4)\in\{+,-\}^4$, set
\[
 r_\sigma=\sum_{i=1}^4 u_i^{\sigma_i}.
\]
The four selected normals have Gram matrix
\begin{equation}\label{eq:vertexgram}
 H=\frac43I_4-\frac13J_4.
\end{equation}
It is positive semidefinite of rank $3$, and its kernel is spanned by
${\bf1}=(1,1,1,1)$.  Consequently $r_\sigma$ is null and is orthogonal
to all four selected normals.

Let $q$ be an interior point of $P$.  Every facet inequality is then
strict, so
\begin{equation}\label{eq:nonzero}
 \langle q,r_\sigma\rangle
   =\sum_i\langle q,u_i^{\sigma_i}\rangle<0.
\end{equation}
In particular $r_\sigma\ne0$.  The four selected normals are linearly
independent: if their linear combination with coefficient vector $c$
vanishes, then $Hc=0$, hence $c=\lambda{\bf1}$; the relation becomes
$\lambda r_\sigma=0$, so $\lambda=0$.

The vertex indexed by $\sigma$ cannot be finite.  Indeed, if a timelike
representative of that vertex existed, the four incident normals would
belong to its positive-definite four-dimensional orthogonal complement,
contrary to the singularity of \eqref{eq:vertexgram}.  Hyperideal
vertices are excluded by hypothesis.  The vertex is therefore ideal.
The common orthogonal complement of its four independent incident
normals is one-dimensional, so its null ray is exactly the ray of
$r_\sigma$; inequality \eqref{eq:nonzero} selects its future-pointing
representative.

\section{The forced Gram matrix}

Put
\[
 x_i=\langle u_i^+,u_i^-\rangle,\qquad
 d_i=u_i^+-u_i^-,\qquad s_i=u_i^++u_i^-.
\]
Fix the other three signs and compare the two vertices obtained by
choosing $+$ or $-$ in coordinate $i$.  Their representatives $r_+$
and $r_-$ are nonproportional future null vectors, because the two
vertices are distinct faces in the assumed cube face lattice.  Two
nonproportional future null vectors have strictly negative Lorentzian
product.  Since $d_i=r_+-r_-$, it follows that
\begin{equation}\label{eq:dpositive}
 d_i^2=-2\langle r_+,r_-\rangle>0.
\end{equation}

Expanding with \eqref{eq:cross} gives
\begin{align}
 \langle d_i,d_j\rangle&=0 &&(i\ne j),\label{eq:dorth}\\
 \langle s_i,d_j\rangle&=0 &&(1\leq i,j\leq4),\label{eq:sdorth}\\
 d_i^2&=2-2x_i,\qquad s_i^2=2+2x_i,\label{eq:squares}\\
 \langle s_i,s_j\rangle&=-\frac43 &&(i\ne j).\label{eq:scross}
\end{align}
By \eqref{eq:dpositive}--\eqref{eq:dorth}, the four vectors $d_i$
span a positive four-space $D$.  Its orthogonal complement in the
space of signature $(4,1)$ is a negative line.  Equation
\eqref{eq:sdorth} places every $s_i$ on this line.  Write
$s_i=a_i\tau$ with $\tau^2=-1$.  Equation \eqref{eq:scross} says
$a_i a_j=4/3$ whenever $i\ne j$.  None of the $a_i$ vanishes, and these
pairwise equations force all four coefficients to be equal and to have
square $4/3$.  Hence $s_i^2=-4/3$, and \eqref{eq:squares} yields
\begin{equation}\label{eq:opposite}
 x_i=-\frac53\qquad(1\leq i\leq4).
\end{equation}

Thus, in the order
$(u_1^+,u_1^-,u_2^+,u_2^-,u_3^+,u_3^-,u_4^+,u_4^-)$, the full Gram
matrix is
\begin{equation}\label{eq:fullgram}
G=\frac13\begin{pmatrix}
 3&-5&-1&-1&-1&-1&-1&-1\\
-5& 3&-1&-1&-1&-1&-1&-1\\
-1&-1& 3&-5&-1&-1&-1&-1\\
-1&-1&-5& 3&-1&-1&-1&-1\\
-1&-1&-1&-1& 3&-5&-1&-1\\
-1&-1&-1&-1&-5& 3&-1&-1\\
-1&-1&-1&-1&-1&-1& 3&-5\\
-1&-1&-1&-1&-1&-1&-5& 3
\end{pmatrix}.
\end{equation}
It has rank $5$ and nonzero inertia $(4,1)$.  Its three-dimensional
kernel is spanned by the relations saying that each of the last three
pair sums equals the first pair sum.  These facts also follow directly
from the decomposition into the positive difference space and the
common negative sum line.

\section{From normals to the oriented polytope}

For the model system define
\begin{equation}\label{eq:modelnormals}
 w_i^\pm=\left(\frac1{\sqrt3},\ \pm\frac2{\sqrt3}e_i\right).
\end{equation}
These eight unit spacelike vectors have Gram matrix
\eqref{eq:fullgram}.  Both ordered normal systems span $V$: their four
difference vectors span a positive four-space, while any pair sum
supplies the negative orthogonal line.

We record why equality of the two Gram matrices yields a well-defined
Lorentz isometry despite the three relations.  For $c\in\mathbb R^8$,
write
\[
 A(c)=\sum c_{i,\pm}u_i^\pm,
 \qquad B(c)=\sum c_{i,\pm}w_i^\pm.
\]
If $A(c)=0$, Gram equality makes $B(c)$ orthogonal to every model
normal.  The model normals span the nondegenerate space $V$, so
$B(c)=0$.  Therefore $T(A(c))=B(c)$ defines a surjective linear
isometry $T\in O(4,1)$ satisfying $T(u_i^\pm)=w_i^\pm$.

This isometry preserves time orientation.  For the interior point $q$
chosen above,
\[
 \langle Tq,w_i^\pm\rangle=\langle q,u_i^\pm\rangle<0.
\]
Adding the two inequalities for any fixed $i$ and using
$w_i^++w_i^-=(2/\sqrt3,0,0,0,0)$ gives $(Tq)_0>0$.  Hence $T$ maps the
future hyperboloid sheet to itself and retains all eight halfspace
orientations.  For every future hyperboloid point $z$,
\[
 \langle z,u_i^\pm\rangle\leq0
 \quad\Longleftrightarrow\quad
 \langle Tz,w_i^\pm\rangle\leq0.
\]
Thus $T$ maps $P$, not merely its supporting hyperplanes, onto the
model halfspace intersection.

In Klein coordinates $y_k=z_k/z_0$, the two inequalities coming from
\eqref{eq:modelnormals} are
\[
 -\frac12\leq y_i\leq\frac12.
\]
The Euclidean cube $[-1/2,1/2]^4$ lies in the closed Klein ball, and
each of its sixteen vertices has squared norm $1$.  Its vertices are
therefore ideal, and its hyperbolic interior is precisely the regular
ideal equiangular cube $C^4_{\alpha _0}$.  This proves
Theorem~\ref{thm:main}.

\section{Exact check, scope, and assistance disclosure}

As an arithmetic check, the rational matrix \eqref{eq:fullgram} was
recomputed in SymPy~1.14.0.  Exact row reduction and exact eigenvalue
calculation give rank $5$, spectrum
\[
 -\frac83\ (1\text{ time}),\qquad 0\ (3\text{ times}),\qquad
 \frac83\ (4\text{ times}),
\]
and verify the three displayed pair-sum relations.  This computation
only checks finite matrix arithmetic; it is not used in place of the
vertex-ray, time-orientation, or halfspace arguments.  The same
distinction between a mathematical claim and the precise scope of its
exact verification is emphasized in a recent discrete-geometry
certificate study \cite{ZengEtAl2026Certificates}; here the deductive
proof above is self-contained.

OpenAI Codex assisted with the bounded literature search, development
and adversarial checking of the Lorentz--Gram argument, the exact
matrix calculation, and preparation of the manuscript.  Its output is
not treated as mathematical evidence.  No proof assistant was used.
The literature comparison is limited to the cited and inspected
sources and is not a claim of priority.  The theorem concerns only
$4$-cube combinatorics at $\alpha_0$; it asserts no volume extremum,
Coxeter classification, right-angled result, different combinatorial
type, or higher-dimensional analogue.

\bibliographystyle{plainurl}
\bibliography{references}

\end{document}
